\section{Introdução}

% Objetivo da seção: apresentar o problema (modelos treinados retêm informação
% sobre titulares; "eliminar" não tem sentido técnico óbvio), enunciar a tese
% da indistinguibilidade e antecipar a estrutura do artigo.

Esta seção deve situar o leitor no descompasso entre o direito à eliminação,
pensado para bancos de dados, e os modelos de aprendizado de máquina, que não
armazenam registros mas parâmetros influenciados por eles. Deve enunciar a tese
central --- eliminar dados de um modelo significa tornar a contribuição do
titular indistinguível da de um modelo que nunca o viu, julgada contra o
atacante que cada lei pressupõe --- e justificar a centralidade da LGPD: os
arts.~16 e 18, IV oferecem encaixe natural para o padrão, e o direito europeu
(RGPD, EDPB, TJUE) entra como contraponto que confirma a relatividade do
conceito de dado pessoal.

\todo{Gancho inicial: caso concreto ou decisão (EDPS v SRB? Parecer 28/2024?).}
\todo{Delimitar escopo: modelos discriminativos e generativos; excluir bancos de dados vetoriais?}
\todo{Enunciar contribuições em lista (3--4 itens).}
\todo{Roteiro das seções.}
